\begin{center}
  {\Large\itshape Alban à l'heure des troubadours -- XIII\textsuperscript{ème} siècle}
\end{center}

\vspace{1.5em}

\noindent
\lettrine{S}{ur} le portail nord de N.D. de Paris, est sculptée dans la
pierre l'histoire de Théophile qui vendit son âme au diable et fut
sauvé par la Vierge. Entre les deux guerres mondiales, un professeur
de Sorbonne et ses étudiants, penchés sur le récit qu'en avait fait
au XII\textsuperscript{ème} siècle le poète Rutebeuf, y trouvèrent le
chemin de la foi. Gustave Cohen passa au catholicisme, un de ses
élèves se fit dominicain, une autre entra au cloître ? Alban n'a pas
eu la même chance. C'est de bouche à oreille qu'est venue jusqu'à
nous une semblable légende, mais combien plus pittoresque que celle
de Paris. Quel troubadour l'inventa, la chanta ? Est-ce un descendant
de ces « Brossard d'Alban », d'origine normande ? L'un d'eux fut le
grand argentier de Philippe le Bel. Veut-il s'implanter dans le midi
comme le fit l'un d'entre eux en 1950 ? Est-ce un ancêtre de ce
« Pierre d'Alban », cabaliste cité par Pierre Benoît dans
« La Châtelaine du Liban » ? Peu importe le nom de ce marchand de
rêve, qui allait jouant de sa viole, chantant les amours de quelque
Tristan et Iseult, ou mieux, célébrant avec ferveur cette Madone de
chez nous, plus suave que le parfum de violette, plus délectable que
le miel le plus savoureux et finalement se laissa prendre à la
« caresse de ces mots ». « Notre Dame »... que de foyers durent
l'entendre.
